    \chapter{Prizes and Congresses in 2026}

    This part records the mathematical events most relevant to the 2026
    Science Quiz season. Its time window begins on January 1, 2025, and runs
    through the preparation cutoff of August 28, 2026. Later developments may be
    added before the competition. Unlike Part II, which collects durable facts
    that should remain useful in future years, this part emphasizes why a
    person, prize, theorem, conjecture, institution, or competition became
    especially important during 2025 or 2026.

    The guiding association is
    \[
        \begin{aligned}
            \text{event}
            &\quad\longleftrightarrow\quad
            \text{person or institution}\\
            &\quad\longleftrightarrow\quad
            \text{mathematical achievement}
            \quad\longleftrightarrow\quad
            \text{year and status}.
        \end{aligned}
    \]

    \begin{fact}[Status language in this part]
        Recent mathematics must be described with care.
        \begin{itemize}
            \item \emph{Published} means that the work appeared in a scholarly
                  publication.
            \item \emph{Accepted} means that a journal announced acceptance, even
                  if the final issue had not yet appeared.
            \item \emph{Preprint} means that a manuscript was publicly posted but
                  had not completed ordinary journal publication.
            \item \emph{Formally verified} means that a proof assistant checked a
                  formal proof term; this is different from ordinary peer review.
            \item \emph{Scheduled} describes an officially announced future event
                  or appointment.
        \end{itemize}
        Numerical evidence, an announced proof, a preprint, a formally verified
        argument, and a published theorem are therefore not interchangeable.
    \end{fact}

    \begin{strategy}
        Past Science Quiz questions frequently combine a recent event with an
        older mathematical fact. For each entry, remember the recent year and
        the shortest unmistakable mathematical anchor. Dates, locations,
        institutions, prize combinations, unusual records, and changes of
        affiliation are especially useful for multiple-choice elimination.
    \end{strategy}

    The chapter on mathematical prizes and congresses in Part II explains the
    Fields Medal, Abel Prize, and Wolf Prize. The present chapter records the
    season-specific results announced in 2025 and 2026.

    \vspace{1em}

    {\small
    \renewcommand{\arraystretch}{1.18}
    \begin{longtable}{@{}L{0.08\textwidth}L{0.25\textwidth}L{0.27\textwidth}L{0.32\textwidth}@{}}
        \toprule
        Year & Prize or congress & Laureate(s) or location & Principal Quiz anchors \\
        \midrule
        \endfirsthead

        \toprule
        Year & Prize or congress & Laureate(s) or location & Principal Quiz anchors \\
        \midrule
        \endhead

        \midrule
        \multicolumn{4}{r}{Continued on the next page.} \\
        \endfoot

        \bottomrule
        \endlastfoot

        2025
        & Abel Prize
        & Masaki Kashiwara
        & Algebraic analysis, representation theory, $D$-modules, crystal bases. \\

        2025
        & Breakthrough Prize in Mathematics
        & Dennis Gaitsgory
        & Geometric Langlands program in characteristic zero. \\

        2025
        & Shaw Prize in Mathematical Sciences
        & Kenji Fukaya
        & Fukaya category, symplectic topology, mirror symmetry, gauge theory. \\

        2025
        & ICTP--IMU Ramanujan Prize
        & Claudio Mu\~noz
        & Nonlinear dispersive PDE, soliton stability, multi-soliton dynamics. \\

        2025
        & Basic Science Lifetime Awards
        & Shigefumi Mori; George Lusztig; Robert Tarjan
        & Minimal model program; representation theory; algorithms and data structures. \\

        2026
        & Abel Prize
        & Gerd Faltings
        & Arithmetic geometry, Mordell conjecture, Lang conjectures, Diophantine geometry. \\

        2026
        & International Congress of Mathematicians
        & Philadelphia, United States, July 23--30
        & Fields Medals and the principal IMU prizes; IMU General Assembly in New York City. \\

        2026
        & Fields Medals
        & Yu Deng; John Pardon; Jacob Tsimerman; Hong Wang
        & PDE and kinetic theory; symplectic geometry; arithmetic geometry and
          o-minimality; harmonic analysis and Kakeya. \\

        2026
        & IMU Abacus Medal
        & Shayan Oveis Gharan
        & Algorithms, combinatorial optimization, approximation algorithms. \\

        2026
        & Carl Friedrich Gauss Prize
        & Yurii Nesterov
        & Convex optimization and accelerated first-order methods. \\

        2026
        & Chern Medal Award
        & Graeme Segal
        & Topology, geometry, and mathematical physics. \\

        2026
        & Leelavati Award
        & Hannah Fry
        & Public communication and popular understanding of mathematics. \\

        2026
        & Breakthrough Prize in Mathematics
        & Frank Merle
        & Nonlinear evolution equations, singularity formation, stability, solitons. \\

        2026
        & Shaw Prize in Mathematical Sciences
        & Emmanuel Cand\`es; Camillo De Lellis
        & Compressed sensing, signal processing, statistics; geometric measure
          theory, singularities, and fluid dynamics. \\

        2026
        & Clay Research Awards
        & Three research groups
        & Furstenberg and Kakeya sets; Telescope Conjecture counterexamples;
          Boltzmann equation from hard-sphere dynamics. \\

        2026
        & Rolf Schock Prize in Mathematics
        & Tobias Holck Colding
        & Minimal surfaces and geometric flows. \\

        2026
        & Maryam Mirzakhani Prize in Mathematics
        & Roman Bezrukavnikov
        & Geometric representation theory and modular representation theory. \\

        2026
        & Michael and Sheila Held Prize
        & Irit Dinur; Subhash Khot; Guy Kindler; Dor Minzer; Muli Safra
        & The $2$-to-$2$ Games Theorem, approximation hardness, Unique Games. \\

        2026
        & International Congress of Basic Science
        & Beijing, China, August 9--21
        & Inaugural ICBS Medals and Frontiers of Science Awards; mathematics,
          physics, and engineering. \\

        2026
        & ICBS Medals in Mathematics
        & Claire Voisin; Horng-Tzer Yau; Shou-Wu Zhang
        & Hodge theory and algebraic geometry; random matrices and mathematical
          physics; arithmetic geometry. \\
    \end{longtable}
    }

    \begin{fact}[ICM 2026]
        The 2026 International Congress of Mathematicians was held in
        Philadelphia from July 23 to July 30. The twentieth IMU General Assembly
        met in New York City immediately before the Congress, on July 20--21.
        The opening ceremony included the Fields Medals, IMU Abacus Medal, Gauss
        Prize, and Chern Medal Award; the Leelavati Award was associated with the
        closing of the Congress.
    \end{fact}

    \begin{fact}[ICBS 2026 and the inaugural mathematics medals]
        The International Congress of Basic Science was held in Beijing from
        August 9 through August 21. The ICBS Medal was introduced
        in 2026 with three medals in mathematics:
        \[
            \begin{aligned}
                \text{Claire Voisin}
                &\quad\longleftrightarrow\quad
                \text{Emmy Noether Medal},\\
                \text{Horng-Tzer Yau}
                &\quad\longleftrightarrow\quad
                \text{Shiing-Shen Chern Medal},\\
                \text{Shou-Wu Zhang}
                &\quad\longleftrightarrow\quad
                \text{Andrew Wiles Medal}.
            \end{aligned}
        \]
        The shortest mathematical anchors are Hodge theory and algebraic
        geometry for Voisin, random matrix theory and mathematical physics for
        Yau, and arithmetic geometry for Zhang. The Congress also presents the
        Frontiers of Science Awards, which recognize research papers rather than
        functioning as another single-laureate mathematics prize.
    \end{fact}

    \vspace{1em}

    {\small
    \renewcommand{\arraystretch}{1.2}
    \begin{longtable}{@{}L{0.20\textwidth}L{0.24\textwidth}L{0.48\textwidth}@{}}
        \toprule
        2026 Fields Medalist & Main area & Season-specific Quiz anchors \\
        \midrule
        \endfirsthead

        \toprule
        2026 Fields Medalist & Main area & Season-specific Quiz anchors \\
        \midrule
        \endhead

        \bottomrule
        \endlastfoot

        Yu Deng
        & Partial differential equations and kinetic theory
        & Rigorous derivation of kinetic equations from many-particle dynamics;
          hard spheres, the Boltzmann equation, wave kinetic equations, and a
          concrete part of Hilbert's sixth-problem program. \\

        John Pardon
        & Symplectic geometry and topology
        & Holomorphic curves, virtual fundamental cycles, Fukaya categories,
          symplectic topology, knot theory, and $3$-manifolds. \\

        Jacob Tsimerman
        & Arithmetic and complex algebraic geometry
        & O-minimality, the Andr\'e--Oort conjecture for Siegel modular varieties,
          period maps, and the Griffiths conjecture. \\

        Hong Wang
        & Harmonic analysis and geometric measure theory
        & Restriction theory, local smoothing, Furstenberg sets, and the
          three-dimensional Kakeya conjecture. She became the third woman to
          receive a Fields Medal. \\
    \end{longtable}
    }

    \begin{fact}[The 2025 and 2026 Abel Prizes]
        The two Abel laureates in this preparation window are distinguished by
        sharply different anchors.
        \[
            \begin{aligned}
                \text{Masaki Kashiwara}
                &\quad\longleftrightarrow\quad
                D\text{-modules and crystal bases},\\
                \text{Gerd Faltings}
                &\quad\longleftrightarrow\quad
                \text{Mordell conjecture and arithmetic geometry}.
            \end{aligned}
        \]
        Faltings had already received the Fields Medal in 1986. Thus he joined
        the group of mathematicians who have received both the Fields Medal and
        the Abel Prize.
    \end{fact}

    \begin{fact}[2026 Clay Research Awards]
        The three recognized clusters are worth remembering as complete
        associations.
        \begin{itemize}
            \item Tuomas Orponen, Pablo Shmerkin, Hong Wang, and Joshua Zahl:
                  Furstenberg sets in the plane and Kakeya sets in
                  three-dimensional space.
            \item Robert Burklund, Jeremy Hahn, Ishan Levy, and Tomer Schlank:
                  counterexamples to the Ravenel Telescope Conjecture in
                  chromatic homotopy theory.
            \item Yu Deng and Zaher Hani: the long-time derivation of the
                  Boltzmann equation from hard-sphere dynamics; Xiao Ma was a
                  coauthor of the mathematical work.
        \end{itemize}
    \end{fact}

    \begin{remark}
        Only officially confirmed prize results are recorded. In particular,
        no unannounced 2026 Wolf Prize in Mathematics laureate is inferred or
        supplied merely to complete the traditional prize triad.
    \end{remark}

    \chapter{Mathematical Competitions in 2026}

    Competition questions differ from prize questions. The most useful chain is
    \[
        \text{competition}
        \quad\longleftrightarrow\quad
        \text{host and year}
        \quad\longleftrightarrow\quad
        \text{team result}
        \quad\longleftrightarrow\quad
        \text{individual record}.
    \]

    \vspace{1em}

    {\small
    \renewcommand{\arraystretch}{1.2}
    \begin{longtable}{@{}L{0.12\textwidth}L{0.21\textwidth}L{0.27\textwidth}L{0.32\textwidth}@{}}
        \toprule
        Competition & Host and scale & Republic of Korea & Notable records \\
        \midrule
        \endfirsthead

        \toprule
        Competition & Host and scale & Republic of Korea & Notable records \\
        \midrule
        \endhead

        \bottomrule
        \endlastfoot

        IMO 2025
        & Sunshine Coast, Australia; $110$ countries and $630$ contestants
        & Third place with $203$ points; four gold and two silver medals
        & Gold cutoff $35$; five perfect scores; six problems over two contest days. \\

        IMO 2026
        & Shanghai, China; $117$ teams and $666$ contestants
        & Sixth place with $167$ points; three gold, two silver, and one
          honorable mention
        & Seven perfect scores; Hyeonjun Lee of the Republic of Korea obtained $42/42$. \\

        APMO 2025
        & Thirty-seventh Asian Pacific Mathematical Olympiad
        & Kyungjun Park scored $35/35$ and received gold; the next nine Korean
          results ranged from silver through honorable mention
        & Five problems, each worth seven points; official results. \\

        APMO 2026
        & Thirty-eighth Asian Pacific Mathematical Olympiad; $43$ countries and
          $397$ contestants
        & Second place with $306$ points; one gold, two silver, four bronze, and
          three honorable mentions
        & Seohyung Jo, Hyeonjun Lee, and Hyunjun Jang each scored $35/35$;
          official awards were gold, silver, and silver, respectively. \\
    \end{longtable}
    }

    \begin{fact}[Republic of Korea at IMO 2025]
        The six Korean contestants and their scores were
        \[
            37,\ 36,\ 36,\ 35,\ 31,\ 28.
        \]
        Kyungjun Park, Wooju Ham, Hyeonjun Lee, and Hyewon Yoon received gold;
        Hyunjun Jang and Hyungjune Cho received silver. The team total was
        \[
            37+36+36+35+31+28=203.
        \]
    \end{fact}

    \begin{fact}[Republic of Korea at IMO 2026]
        The Republic of Korea placed sixth with $167$ points. Its aggregate problem scores were
        \[
            42,\ 20,\ 17,\ 40,\ 37,\ 11.
        \]
        Hyeonjun Lee was one of seven contestants worldwide to obtain a perfect
        score of $42$.
    \end{fact}

    \begin{fact}[Republic of Korea at APMO 2026]
        The official APMO 2026 results place the Republic of Korea second among
        $43$ participating countries, with a team total of $306$. Seohyung Jo,
        Hyeonjun Lee, and Hyunjun Jang each obtained the perfect score $35/35$.
        The official award labels were
        \[
            \text{Seohyung Jo: Gold},
            \qquad
            \text{Hyeonjun Lee: Silver},
            \qquad
            \text{Hyunjun Jang: Silver}.
        \]
        Thus equal perfect scores can receive different official award labels in
        APMO 2026; the numerical score and the award label should be remembered
        separately.
    \end{fact}

    \begin{fact}[The 2026 IMO medal design]
        The reverse of the 2026 IMO medal used Zhao Shuang's diagram for the
        Pythagorean theorem, a classical diagram associated with his commentary
        on the \textit{Zhoubi Suanjing}. The medal had diameter $9\,\mathrm{cm}$.
        This provides a direct connection between a current competition, the
        Pythagorean theorem, and the history of Chinese mathematics.
    \end{fact}

    \begin{remark}
        Human competition results belong in this chapter. Performance by AI
        systems on the same or similar problems is discussed separately, because
        an official contestant result, an external evaluation, a natural-language
        proof, and a machine-checked formal proof are different objects.
    \end{remark}

    \chapter{Major Mathematical Developments in 2026}

    This chapter records developments that were prominent during 2025 and 2026.
    Complete proofs are generally far beyond the scope of Science Quiz. The
    intended memory unit is
    \[
        \text{standard problem or theorem}
        \quad+\quad
        \text{people}
        \quad+\quad
        \text{status}
        \quad+\quad
        \text{main method or field}.
    \]

    \vspace{1em}

    {\small
    \renewcommand{\arraystretch}{1.18}
    \begin{longtable}{@{}L{0.25\textwidth}L{0.22\textwidth}L{0.18\textwidth}L{0.27\textwidth}@{}}
        \toprule
        Development & Principal names & Status in this season & Quiz anchors \\
        \midrule
        \endfirsthead

        \toprule
        Development & Principal names & Status in this season & Quiz anchors \\
        \midrule
        \endhead

        \midrule
        \multicolumn{4}{r}{Continued on the next page.} \\
        \endfoot

        \bottomrule
        \endlastfoot

        Three-dimensional Kakeya conjecture
        & Hong Wang; Joshua Zahl
        & 2025 proof preprint; recognized by a 2026 Clay Research Award
        & Full Hausdorff and Minkowski dimension $3$; harmonic analysis and
          geometric measure theory. \\

        Hilbert's tenth problem over rings of integers
        & Levent Alp\"oge; Manjul Bhargava; Wei Ho; Ari Shnidman; independently
          Peter Koymans and Carlo Pagano
        & One proof accepted and published; an independent preprint
        & Diophantine equations, undecidability, number fields, elliptic curves. \\

        Boltzmann equation from hard spheres
        & Yu Deng; Zaher Hani; Xiao Ma
        & Long-time derivation; recognized by a 2026 Clay Research Award
        & Kinetic theory, many-particle dynamics, Hilbert's sixth problem. \\

        Modularity of ordinary abelian surfaces
        & George Boxer; Frank Calegari; Toby Gee; Vincent Pilloni
        & 2025 preprint
        & Abelian surfaces, Siegel modular forms, Galois representations,
          Langlands philosophy. \\

        Mizohata--Takeuchi conjecture counterexample
        & Hannah Mira Cairo
        & 2025 preprint
        & Harmonic analysis, Fourier restriction, counterexample, logarithmic loss. \\

        Compact Bonnet pairs
        & Alexander Bobenko; Tim Hoffmann; Andrew Sageman-Furnas
        & Published in 2025
        & Noncongruent immersed tori with the same metric and mean curvature. \\

        Optimal spectral gap for typical hyperbolic surfaces
        & Nalini Anantharaman; Laura Monk
        & 2025 preprint completing a multi-paper program
        & Random hyperbolic surfaces, Weil--Petersson measure, Laplacian,
          asymptotically optimal spectral gap $1/4$. \\

        A convex polyhedron without Rupert's property
        & Jakob Steininger; Sergey Yurkevich
        & 2025 preprint
        & The Noperthedron; first proved convex polyhedron that cannot pass
          through a suitable straight hole in an identical copy of itself. \\

        Unknotting number is not additive
        & Mark Brittenham; Susan Hermiller
        & Accepted in 2026
        & Knot theory, connected sum, explicit counterexamples. \\

        Kissing-number lower bounds in dimensions $17$--$21$
        & Henry Cohn; Anqi Li
        & Preprint widely publicized in 2025
        & Sphere packing, spherical codes, nonlattice configurations. \\

        Furstenberg set conjecture
        & Tuomas Orponen; Pablo Shmerkin and collaborators
        & Major planar result; 2026 Clay recognition
        & Geometric measure theory, sets containing line segments in many directions. \\

        Ravenel Telescope Conjecture
        & Robert Burklund; Jeremy Hahn; Ishan Levy; Tomer Schlank
        & Counterexamples recognized in 2026
        & Chromatic homotopy theory and stable homotopy theory. \\

        Jacobian Conjecture
        & Levent Alpöge; AI-assisted discovery involving Claude Fable
        & Explicit $3$-dimensional counterexample announced and independently
          formally verified in 2026; counterexamples now exist for every $n>2$
        & Polynomial maps, constant Jacobian determinant, noninjectivity; the
          case $n=2$ remains open. \\

        Riemann zeta zeros on the critical line
        & Levent Alp\"oge; Ralph Furman; mathematical argument discovered and
          written by Claude
        & August 2026 preprint; Lean formalization accompanies the paper
        & Unconditionally more than two thirds of the nontrivial zeros are
          simple and on $\Real(s)=1/2$; the Riemann Hypothesis remains open. \\

        Sendov's Conjecture
        & Lech Mazur; Terence Tao
        & AI-assisted proof in August 2026; human-digested exposition and Lean
          formalization
        & Zeros of a polynomial in the unit disk and nearby critical points;
          Phelps--Rodriguez strengthening. \\

        Banach's isometric conjecture
        & Xinbao Lu; Kaiwen Yang; Antonio Acuaviva; Tomasz Kania
        & August 2026 preprints completing the real and complex cases and a
          quaternionic counterpart
        & Normed spaces, mutually isometric finite-dimensional subspaces,
          Hilbert and Hermitian inner-product norms. \\

        Gromov's codimension-two volume-growth conjecture
        & Jian Ge; independently Bochao Kong and Xingyu Zhu
        & Independent August 2026 preprints
        & Nonnegative Ricci curvature, positive scalar curvature, and the sharp
          growth order $R^{n-2}$. \\

        Samuels' Conjecture
        & Zhi Ling
        & August 2026 preprint
        & Sharp lower bound for sums of independent nonnegative random variables;
          Feige's conjecture follows from the equal-means case. \\

        Erd\H{o}s--Szemer\'edi sum--product conjecture over $\R$
        & Thomas F. Bloom; Will Sawin; Carl Schildkraut; Dmitrii Zhelezov
        & Disproved by a May 2026 preprint
        & Additive combinatorics, sumsets and product sets, algebraic number
          fields; the integer version is not disproved. \\

        Cycle Double Cover Conjecture
        & OpenAI; expositions by Sang-il Oum and Jim Geelen
        & Affirmative AI-generated proof announced in July 2026 and publicly
          exposited in preprints
        & Bridgeless graphs, cycle covers, every edge covered exactly twice. \\

        Maxwell Conjecture on point charges
        & Philip Arathoon; Gavin Ball; Matthew D. Kvalheim
        & Disproved by a July 2026 preprint
        & Five point charges with at least $24$ nondegenerate critical points,
          exceeding Maxwell's proposed $(n-1)^2$ bound. \\

        Planar Pompeiu, Schiffer, and Berenstein conjectures
        & Matthew J. Colbrook; George Stepaniants; Siavash Sadeghi
        & Computer-assisted counterexample preprints in August 2026
        & Overdetermined eigenvalue problems, Fourier zero sets, noncircular
          real-analytic planar domains, interval arithmetic. \\

        Petersen Coloring Conjecture
        & Bryce Putman; Jorik Jooken
        & Explicit $112$-vertex counterexample posted in August 2026; SAT/DRAT
          verification followed by a human-checkable proof
        & Bridgeless cubic graphs, Petersen coloring, normal $5$-edge-coloring. \\

        Heil--Ramanathan--Topiwala Conjecture
        & Markus Faulhuber; Philipp Petersen; Jordy Timo van Velthoven; Felix
          Voigtlaender; Vignon Oussa
        & Disproved in August 2026; a $12$-shift counterexample was followed by
          a computer-assisted $4$-shift counterexample
        & Time--frequency analysis, Gabor systems, linear dependence of
          time--frequency shifts of a Schwartz function. \\

    \end{longtable}
    }

    \begin{theorem}[Three-Dimensional Kakeya Conjecture]
        A Kakeya set in $\R^3$ contains a unit line segment in every direction.
        The 2025 result of Hong Wang and Joshua Zahl establishes that every such
        set has Hausdorff dimension and Minkowski dimension equal to $3$.
    \end{theorem}

    \begin{idea}
        The theorem does not say that every Kakeya set has positive
        three-dimensional volume. It says that no Kakeya set can have dimension
        smaller than the ambient dimension. The main quiz association is
        \[
            \text{Kakeya in }\R^3
            \quad\longleftrightarrow\quad
            \text{Hong Wang and Joshua Zahl}.
        \]
        Both the Hausdorff dimension and the Minkowski dimension are equal to
        $3$.
    \end{idea}

    \begin{fact}[Hilbert's tenth problem over number fields]
        Hilbert's tenth problem asks for an algorithm deciding whether an
        integer polynomial equation has an integer solution. Matiyasevich,
        building on Davis, Putnam, and Robinson, proved that no such general
        algorithm exists over $\Z$.

        A paper by Levent Alp\"oge, Manjul Bhargava, Wei Ho, and Ari
        Shnidman was accepted and published during the 2025--2026 season and
        established the analogous negative answer for the ring of integers of
        every number field. Peter Koymans and Carlo Pagano posted an independent
        and more general undecidability result for infinite rings finitely
        generated over $\Z$.
    \end{fact}

    \begin{remark}
        The phrase ``Hilbert's tenth problem over every number field'' refers to
        the corresponding ring of algebraic integers. It is not the same as the
        still distinct question of Hilbert's tenth problem over the rational
        numbers $\Q$.
    \end{remark}

    \begin{fact}[Boltzmann equation and Hilbert's sixth problem]
        Yu Deng, Zaher Hani, and Xiao Ma developed a rigorous long-time
        derivation of the Boltzmann equation from deterministic hard-sphere
        dynamics. This is a major realization of the kinetic-theory part of
        Hilbert's sixth-problem program.
    \end{fact}

    \begin{remark}
        Hilbert's sixth problem is a broad request for an axiomatic treatment of
        physics. The hard-sphere-to-Boltzmann result is therefore a major
        solution of a concrete component, not a declaration that every aspect of
        Hilbert's sixth problem has been completed.
    \end{remark}

    \begin{fact}[Modularity beyond elliptic curves]
        George Boxer, Frank Calegari, Toby Gee, and Vincent Pilloni proved
        modularity for an important class of ordinary abelian surfaces over
        $\Q$, under explicit local and residual hypotheses. In particular, their
        work proves modularity for a positive proportion of abelian surfaces in
        a natural ordering.
    \end{fact}

    \begin{idea}
        An elliptic curve is one-dimensional as an abelian variety, whereas an
        abelian surface has dimension $2$. The association to retain is
        \[
            \text{elliptic curves and modular forms}
            \quad\longrightarrow\quad
            \text{abelian surfaces and Siegel modular forms}.
        \]
        This is part of the broader Langlands philosophy connecting arithmetic
        objects with automorphic forms and representations.
    \end{idea}

    \begin{fact}[Hannah Cairo's counterexample]
        In 2025, Hannah Mira Cairo constructed a counterexample to the
        Mizohata--Takeuchi conjecture in Fourier restriction theory. Her result
        showed that the conjectured estimate fails in its original form and
        identified a logarithmic loss that cannot simply be removed.
    \end{fact}

    \begin{fact}[Compact Bonnet pairs]
        A Bonnet pair consists of noncongruent surfaces with the same induced
        metric and the same mean-curvature function. In 2025, Alexander Bobenko,
        Tim Hoffmann, and Andrew Sageman-Furnas published the first compact
        examples: a pair of immersed tori.
    \end{fact}

    \begin{fact}[Unknotting number is not additive]
        Let $u(K)$ denote the unknotting number of a knot $K$. The long-standing
        additivity expectation was
        \[
            u(K_1\mathbin{\#}K_2)
            =
            u(K_1)+u(K_2).
        \]
        Brittenham and Hermiller constructed examples for which
        \[
            u(K_1\mathbin{\#}K_2)
            <
            u(K_1)+u(K_2),
        \]
        and their paper was accepted in 2026. Thus a connected sum may require
        fewer crossing changes than separately unknotting the two summands.
    \end{fact}

    \begin{fact}[Improved kissing configurations]
        Henry Cohn and Anqi Li improved the best known lower bounds for kissing
        numbers in dimensions $17$ through $21$. Their constructions used
        spherical codes that need not come from lattices, reinforcing the lesson
        that the most symmetric construction need not be the best known one.
    \end{fact}

    \begin{fact}[More than two thirds of the zeta zeros are simple and on the critical line]
        In an August 2026 preprint, Levent Alp\"oge and Ralph Furman proved
        unconditionally that at least two thirds of the nontrivial zeros of the
        Riemann zeta function, counted with multiplicity, are simple and lie on
        the critical line
        \[
            \Real(s)=\frac12.
        \]
        They also proved that at least five sixths are distinct. With the
        Montgomery--Taylor window, these lower bounds improve to approximately
        \[
            0.6725
            \qquad\text{and}\qquad
            0.8362.
        \]
        The proof replaces a positivity step that classically used the Riemann
        Hypothesis with a rank--trace inequality for a finite compression of
        Weil's Hermitian form, together with Sylvester's law of inertia. The
        paper includes a Lean~4 formalization.
    \end{fact}

    \begin{remark}
        This is a new unconditional record toward the Riemann Hypothesis, not a
        proof of the hypothesis. The paper states that the mathematical argument
        was discovered and written by Claude, with Alp\"oge and Furman verifying
        the proof and taking responsibility for its contents.
    \end{remark}

    \begin{fact}[Sendov's conjecture and the Phelps--Rodriguez strengthening]
        Sendov's conjecture concerns a polynomial whose zeros all lie in the
        closed unit disk: every zero must lie within distance $1$ of some
        critical point. In August 2026, Lech Mazur obtained an AI-assisted proof
        for all degrees. Terence Tao then produced a human-digested exposition
        and a substantially shorter Lean formalization, and observed that the
        same argument proves the stronger Phelps--Rodriguez statement. The
        durable statement and its 2026 status are also recorded in Facts to
        Memorize.
    \end{fact}

    \begin{fact}[Banach's isometric conjecture]
        Banach asked in 1932 whether a real normed space must be a Hilbert space
        if, for some fixed dimension $n>1$, all of its $n$-dimensional subspaces
        are mutually isometric. Gromov had proved the even-dimensional case.
        In an August 2026 preprint, Xinbao Lu and Kaiwen Yang proved every
        remaining odd-dimensional case, completing the conjecture over the real
        field. A follow-up preprint by Antonio Acuaviva and Tomasz Kania completed
        the complex version, where the norm is forced by a Hermitian inner
        product, and also established the quaternionic counterpart.
    \end{fact}

    \begin{fact}[Gromov's codimension-two volume-growth conjecture]
        In 1986, Gromov asked whether a complete noncompact $n$-dimensional
        Riemannian manifold with nonnegative Ricci curvature and scalar curvature
        bounded below by $1$ must satisfy
        \[
            \vol(B(p,R))\leq C_nR^{n-2}
        \]
        for all centers $p$ and radii $R>0$. Independent August 2026 preprints by
        Jian Ge and by Bochao Kong with Xingyu Zhu gave affirmative resolutions.
        The exponent $n-2$ is the reason for the name ``codimension-two volume
        growth.''
    \end{fact}

    \begin{fact}[Samuels' conjecture and Feige's conjecture]
        Samuels' conjecture asks for the sharp smallest possible probability that
        a sum of independent nonnegative random variables stays below a given
        threshold when their means are prescribed. In an August 2026 preprint,
        Zhi Ling proved the conjectured sharp bound and showed that it is
        attained. Feige's conjecture follows immediately from the equal-means
        case, giving the most useful short association
        \[
            \text{Samuels}\quad\longrightarrow\quad
            \text{Feige as the equal-means case}.
        \]
    \end{fact}

    \begin{fact}[Sum--product conjecture over the reals]
        For a finite set $A$, write
        \[
            A+A=\{a+b:a,b\in A\},
            \qquad
            AA=\{ab:a,b\in A\}.
        \]
        In May 2026, Thomas F. Bloom, Will Sawin, Carl Schildkraut, and Dmitrii
        Zhelezov constructed arbitrarily large finite $A\subset\R$ satisfying
        \[
            \max\{|A+A|,|AA|\}\leq |A|^{2-c}
        \]
        for an absolute constant $c>0$. This disproves the Erd\H{o}s--Szemer\'edi
        sum--product conjecture over $\R$. Because their sets live in algebraic
        number fields of degree growing like $\log|A|$, the result does not by
        itself settle the original integer case.
    \end{fact}

    \begin{fact}[Maxwell's proposed bound for electrostatic equilibria is false]
        Maxwell proposed that the field of $n$ point charges should have at most
        $(n-1)^2$ nondegenerate critical points. Philip Arathoon, Gavin Ball, and
        Matthew D. Kvalheim gave a configuration of five point charges whose
        electrostatic potential has at least $24$ nondegenerate critical points.
        Since $(5-1)^2=16$, the proposed bound is false.
    \end{fact}

    \begin{fact}[Planar Pompeiu--Schiffer and Berenstein counterexamples]
        In August 2026, Matthew J. Colbrook and George Stepaniants constructed a
        bounded simply connected noncircular planar domain with real-analytic
        boundary that is simultaneously a counterexample to the planar Schiffer
        conjecture and the corresponding Pompeiu conjecture. Colbrook, Siavash
        Sadeghi, and Stepaniants then adapted the method to disprove the
        unrestricted planar Berenstein conjecture. Both results are
        computer-assisted preprints with rigorous interval-arithmetic
        certification.
    \end{fact}

    \begin{fact}[Petersen coloring counterexample]
        Bryce Putman's August 2026 preprint gives an explicit simple bridgeless
        cubic graph on $112$ vertices with no Petersen coloring and no normal
        $5$-edge-coloring. The computational proof supplies SAT instances and
        checked DRAT certificates. Jorik Jooken subsequently posted a
        human-checkable proof for the $112$-vertex counterexample.
    \end{fact}

    \begin{fact}[The HRT conjecture is false]
        The Heil--Ramanathan--Topiwala conjecture asserted that every finite set
        of distinct time--frequency shifts of a nonzero function in
        $L^2(\R)$ is linearly independent. In August 2026, Markus Faulhuber,
        Philipp Petersen, Jordy Timo van Velthoven, and Felix Voigtlaender
        disproved the conjecture using $12$ time--frequency shifts of a
        Schwartz function. Vignon Oussa then produced a computer-assisted
        counterexample using only $4$ shifts, again with a nonzero Schwartz
        function.
    \end{fact}

    \begin{fact}[Two memorable geometry results from 2025]
        Two visually memorable results complement the major theorem-and-prize
        stories of the season.
        \begin{itemize}
            \item Nalini Anantharaman and Laura Monk completed a program showing
                  that a Weil--Petersson random hyperbolic surface of large genus
                  has spectral gap arbitrarily close to the optimal value
                  $1/4$ with probability tending to $1$.
            \item Jakob Steininger and Sergey Yurkevich constructed the
                  \emph{Noperthedron}, a convex polyhedron without Rupert's
                  property. Thus not every convex polyhedron can pass through a
                  suitable straight hole in an identical copy of itself.
        \end{itemize}
    \end{fact}

    \begin{strategy}
        For a newly announced result, do not memorize only the headline
        ``problem solved.'' Memorize whether the event was a proof, a
        counterexample, an improved bound, a construction, or a new framework.
        These distinctions are common sources of plausible but incorrect answer
        choices.
    \end{strategy}

    \chapter{Artificial Intelligence and Mathematics in 2026}

    The years 2025 and 2026 marked a transition from isolated demonstrations to
    sustained use of AI in olympiad reasoning, algorithm discovery, formal proof,
    and mathematical research. The correct comparison is not merely
    ``human versus machine.'' One must identify the type of mathematical output
    and the manner in which it was checked.

    \vspace{1em}

    {\small
    \renewcommand{\arraystretch}{1.2}
    \begin{longtable}{@{}L{0.23\textwidth}L{0.29\textwidth}L{0.40\textwidth}@{}}
        \toprule
        Type of claim & What was produced & What the label does and does not mean \\
        \midrule
        \endfirsthead

        \toprule
        Type of claim & What was produced & What the label does and does not mean \\
        \midrule
        \endhead

        \bottomrule
        \endlastfoot

        Gold-medal-level score
        & A score at or above the human gold-medal cutoff on a specified problem set
        & It is a benchmark statement, not automatically a formal proof or a new
          research theorem. \\

        Natural-language proof
        & A human-readable written argument
        & It may be judged by expert graders, but it is not mechanically checked
          merely because a model produced it. \\

        Formal proof
        & A proof term accepted by a proof assistant such as Lean
        & The kernel checks the formal derivation; the formal statement and its
          relation to the intended theorem must still be correct. \\

        AI-assisted discovery
        & A construction, algorithm, conjecture, counterexample candidate, or proof idea
        & Human mathematical validation, literature search, attribution, and
          publication status remain essential. \\

        Numerical candidate
        & Computational evidence suggesting a phenomenon or singularity
        & It is not a theorem until an appropriate rigorous argument is supplied. \\
    \end{longtable}
    }

    \begin{fact}[Gold-medal-level performance at IMO 2025]
        Google DeepMind reported that an advanced version of Gemini with Deep
        Think solved five of the six IMO 2025 problems and received $35$ out of
        $42$ points under an official evaluation process. Since the human gold
        cutoff was also $35$, the result met the gold-medal standard. The system
        worked from the natural-language problem statements and produced
        natural-language proofs.
    \end{fact}

    \begin{remark}
        The phrase ``gold-medal-level'' is exact but limited: an AI system was not
        an official contestant and did not receive an IMO medal. The statement
        compares its evaluated score with the human cutoff.
    \end{remark}

    \begin{fact}[Two AI systems at IMO 2025]
        OpenAI separately reported that a general-purpose reasoning model also
        scored $35$ out of $42$ on the IMO 2025 problem set. Thus both OpenAI and
        Google DeepMind announced gold-medal-level natural-language performance
        in 2025, using different systems. Neither system was an official
        contestant, and neither received an IMO medal.
    \end{fact}

    \begin{fact}[AlphaProof and formal mathematics]
        A paper on AlphaProof appeared in \textit{Nature} in 2025. AlphaProof
        combines reinforcement learning with the Lean proof assistant and
        produces machine-checkable formal proofs. Its well-known benchmark used
        the 2024 IMO problem set and reached silver-medal-level performance when
        combined with a geometry system.
    \end{fact}

    \begin{fact}[AlphaEvolve]
        Google DeepMind introduced AlphaEvolve in 2025 as an evolutionary coding
        agent for discovering and improving algorithms. Its mathematical
        significance lies in searching large spaces of executable constructions
        and evaluating candidates automatically, rather than merely generating
        prose solutions.
    \end{fact}

    \begin{fact}[ICPC-level algorithmic performance]
        In 2025, Gemini Deep Think was also evaluated on the ICPC World Finals
        problem set and achieved a gold-medal-level result. This is primarily an
        achievement in algorithm design and competitive programming, whereas the
        IMO emphasizes proof-based olympiad mathematics.
    \end{fact}

    \begin{fact}[GPT-5.5 and off-diagonal Ramsey numbers]
        In April 2026, OpenAI reported that an internal GPT-5.5 system with a
        custom research harness found a new proof of a longstanding asymptotic
        fact about off-diagonal Ramsey numbers. The argument was subsequently
        verified in Lean. For quiz purposes, the important association is
        \[
            \text{GPT-5.5}
            \quad\longleftrightarrow\quad
            \text{off-diagonal Ramsey numbers}
            \quad\longleftrightarrow\quad
            \text{Lean verification}.
        \]
    \end{fact}

    \begin{fact}[GPT-5.6 Sol and the Cycle Double Cover Conjecture]
        The Cycle Double Cover Conjecture asks whether every finite bridgeless
        loopless multigraph admits a multiset of cycles in which every edge
        appears exactly twice. In July 2026, OpenAI released the prompt and
        materials for an affirmative proof generated by GPT-5.6 Sol Ultra.
        Sang-il Oum and Jim Geelen posted independent expository write-ups of the
        argument. The season-specific association is therefore
        \[
            \begin{aligned}
                \text{Cycle Double Cover Conjecture}
                &\longleftrightarrow \text{GPT-5.6 Sol Ultra}\\
                &\longleftrightarrow \text{Oum and Geelen expositions}.
            \end{aligned}
        \]
        The public proof and expositions should still be distinguished from the
        separate question of ordinary journal publication.
    \end{fact}

    \begin{fact}[First Proof]
        In February 2026, OpenAI released proof attempts by an internal model for
        all ten problems in the research-level First Proof challenge. After
        expert feedback, OpenAI assessed at least five attempts---Problems $4$,
        $5$, $6$, $9$, and $10$---as having a high chance of being correct, while
        other attempts remained under review. Its initially favored attempt for
        Problem $2$ was later acknowledged to be incorrect. The episode is a
        useful reminder that a persuasive research proof attempt still requires
        independent mathematical checking.
    \end{fact}

    \begin{fact}[AI disproof of the Erd\H{o}s unit-distance conjecture]
        Let $u(n)$ be the largest possible number of pairs at unit distance among
        $n$ points in the plane. Erd\H{o}s had conjectured that
        \[
            u(n)=n^{1+o(1)}.
        \]
        In May 2026, OpenAI released a proof generated by an internal
        general-purpose model that disproved this conjecture. The construction
        gives, for infinitely many $n$ and some fixed $\delta>0$,
        \[
            u(n)\geq n^{1+\delta}.
        \]
        The argument uses class-field towers and other ideas from algebraic
        number theory, and external mathematicians checked the released proof.
        The result disproves the conjectured near-linear upper bound; it does not
        determine the exact asymptotic order of $u(n)$.
    \end{fact}

    \begin{fact}[Ten AI-generated research advances]
        On August 1, 2026, OpenAI released manuscripts for ten results generated
        by an internal version of its Astra model. Humans prepared the
        manuscripts with the same model, after which the model formalized each
        argument as a Lean certificate. The announced results were:
        \begin{bookitems}
            \item new upper bounds for high-dimensional sphere packing down to
                  the Cohn--Elkies threshold;
            \item exponentially improved bounds for binary and spherical codes;
            \item a construction establishing the existence of non-sofic groups;
            \item a disproof of Connes's rigidity conjecture;
            \item new arithmetic-circuit lower bounds for the permanent,
                  including an arithmetic-formula lower bound of order
                  $n^4/\log n$;
            \item an exponential parallel-repetition theorem for general
                  two-player quantum games;
            \item polynomial-factor hardness of approximation for the closest
                  vector problem;
            \item a resolution of Ehrhart's volume conjecture;
            \item a superexponential lower bound for multicolor triangle Ramsey
                  numbers, resolving Erd\H{o}s Problem $183$;
            \item results on compactness and degeneracy in extremal graph theory,
                  resolving Erd\H{o}s Problems $146$ and $180$.
        \end{bookitems}
        At the August 28, 2026 preparation cutoff, these remained recently
        released manuscripts with formal certificates. Formal verification of
        the stated Lean theorems and
        long-term independent assessment by the mathematical community are
        distinct stages.
    \end{fact}

    \begin{fact}[Palomar registry for Lean-verified mathematics]
        On August 18, 2026, a group including Jeremy Avigad, Terence Tao, Ravi
        Vakil, and Akshay Venkatesh announced Palomar, a public registry of
        Lean-verified mathematics incubated by ICARM and the Lean Focused
        Research Organization. Palomar records fixed source versions, exact
        formal statements, dependencies, and review comments. Its central
        methodological point is equally important: successful kernel checking
        verifies a formal derivation from stated assumptions, but by itself does
        not certify novelty, mathematical importance, ordinary peer review, or
        that an informal claim has been formalized exactly as intended.
    \end{fact}

    \begin{fact}[Claude and the Riemann zeta critical-line record]
        An August 2026 paper by Levent Alp\"oge and Ralph Furman states that its
        mathematical argument was discovered and written by Claude, an AI system
        developed by Anthropic; the listed authors verified the proof and take
        responsibility for it. The mathematical result itself is recorded in
        the Major Mathematical Developments chapter: an unconditional lower
        bound of more than two thirds for nontrivial zeta zeros that are both
        simple and on the critical line, accompanied by Lean~4 formalization.
        The Riemann Hypothesis itself remains open.
    \end{fact}

    \begin{fact}[Leiden Declaration on AI and Mathematics]
        The Leiden Declaration, released in 2026 and endorsed by the IMU,
        articulated principles for responsible AI-assisted mathematics. The
        central requirements are transparent disclosure of significant AI use,
        accurate attribution, openness about methods and evidence, and continued
        human responsibility for mathematical correctness and scholarship.
    \end{fact}

    \begin{strategy}
        A Science Quiz question may deliberately mix the following terms:
        \[
            \text{IMO score},\quad
            \text{formal verification},\quad
            \text{peer review},\quad
            \text{AI-assisted discovery},\quad
            \text{new theorem}.
        \]
        Treat each as a separate claim. A result may satisfy more than one, but
        none implies all the others.
    \end{strategy}

    \chapter{Mathematicians and Institutions in 2026}

    Recent-person questions often use a move, retirement, death, or unusual
    career event as the first clue, followed by an older theorem or prize. The
    following tables record the principal institutional changes and memorial
    events of the season.

    \vspace{1em}

    {\small
    \renewcommand{\arraystretch}{1.18}
    \begin{longtable}{@{}L{0.20\textwidth}L{0.19\textwidth}L{0.24\textwidth}L{0.29\textwidth}@{}}
        \toprule
        Mathematician & 2025--2026 event & Institution & Quiz anchors \\
        \midrule
        \endfirsthead

        \toprule
        Mathematician & 2025--2026 event & Institution & Quiz anchors \\
        \midrule
        \endhead

        \bottomrule
        \endlastfoot

        Hong Wang
        & Became a permanent professor in 2025
        & Institut des Hautes \`Etudes Scientifiques, jointly with New York University
        & Harmonic analysis, Kakeya, 2026 Fields Medal. \\

        Jacob Tsimerman
        & Announced in July 2026 that he would begin a position at OpenAI
        & OpenAI; on leave from the University of Toronto
        & Arithmetic geometry, o-minimality, 2026 Fields Medal, AI safety. \\

        Tim Roughgarden
        & Joined the faculty on July 1, 2026
        & Institute for Advanced Study
        & Algorithms, algorithmic game theory, mechanism design. \\

        James Maynard
        & Appointed Regius Professor; scheduled to begin October 1, 2026
        & \UniversityOfOxford
        & Prime gaps, sieve methods, 2022 Fields Medal; successor to Andrew Wiles. \\

        Ryan Chen; Alex Cohen; Anna Skorobogatova
        & Began Clay Research Fellowships in 2025
        & Clay Mathematics Institute appointments hosted at research institutions
        & Clay Research Fellowships and early-career research appointments. \\

        Oliver Edtmair; Qiuyu Ren
        & Began Clay Research Fellowships in 2026
        & Clay Mathematics Institute appointments hosted at research institutions
        & Clay Research Fellowships and early-career research appointments. \\
    \end{longtable}
    }

    \begin{fact}[Jacob Tsimerman and OpenAI]
        After receiving a 2026 Fields Medal, Jacob Tsimerman announced at ICM
        2026 that he would begin a position at OpenAI. He went on leave from the
        University of Toronto to work on AI safety, while remaining a professor
        and faculty member there. The accurate season-specific association is
        therefore a Fields Medalist moving temporarily from ordinary university
        duties into AI-safety research, not an unqualified claim that he quit
        mathematics or resigned from the university.
    \end{fact}

    \vspace{1em}

    {\small
    \renewcommand{\arraystretch}{1.18}
    \begin{longtable}{@{}L{0.21\textwidth}L{0.18\textwidth}L{0.45\textwidth}@{}}
        \toprule
        Mathematician & Death & Principal associations \\
        \midrule
        \endfirsthead

        \toprule
        Mathematician & Death & Principal associations \\
        \midrule
        \endhead

        \bottomrule
        \endlastfoot

        Peter D. Lax
        & May 16, 2025, aged $99$
        & Partial differential equations, conservation laws, shock waves,
          numerical analysis, Lax equivalence theorem, Wolf Prize 1987, Abel Prize 2005. \\

        Heisuke Hironaka
        & March 18, 2026, aged $94$
        & Resolution of singularities in characteristic zero, algebraic geometry,
          Fields Medal 1970. \\

        Michael O. Rabin
        & April 14, 2026
        & Finite automata, probabilistic algorithms, Rabin--Karp string search,
          Turing Award 1976 shared with Dana Scott. \\

        Cleve Moler
        & May 20, 2026, aged $86$
        & Numerical linear algebra and numerical computing; creator of MATLAB;
          cofounder of MathWorks. \\
    \end{longtable}
    }

    \begin{fact}[Peter D. Lax]
        \par\noindent
        Peter D. Lax spent his academic career at New York University's Courant
        Institute. His work connected partial differential equations, numerical
        methods, fluid dynamics, shock waves, scattering theory, and integrable
        systems. The Lax equivalence theorem states, in its standard numerical
        analysis setting, that a consistent finite-difference approximation of
        a well-posed linear initial-value problem is convergent if and only if it
        is stable.
    \end{fact}

    \begin{fact}[Heisuke Hironaka]
        Hironaka received the 1970 Fields Medal for proving resolution of
        singularities for algebraic varieties over fields of characteristic
        zero. The season-specific clue is his death in 2026; the durable
        mathematical clue is
        \[
            \text{Hironaka}
            \quad\longleftrightarrow\quad
            \text{resolution of singularities}.
        \]
    \end{fact}

    \begin{fact}[Cleve Moler and MATLAB]
        Cleve Moler was a pioneer of numerical computing and the creator of
        MATLAB, originally developed to give students convenient interactive
        access to matrix software. He later cofounded MathWorks. His death in
        May 2026 creates a particularly short quiz association:
        \[
            \text{Cleve Moler}
            \quad\longleftrightarrow\quad
            \text{MATLAB and numerical computing}.
        \]
    \end{fact}

    \begin{fact}[From Wiles to Maynard at Oxford]
        Andrew Wiles was the inaugural holder of Oxford's modern Regius
        Professorship of Mathematics. James Maynard was appointed as his
        successor, with the new appointment scheduled to begin on October 1,
        2026. This creates a memorable association between Fermat's Last Theorem
        and modern work on prime gaps.
    \end{fact}

    \chapter{Mathematics in the Republic of Korea in 2026}

    This chapter records events whose specifically Korean institutional or
    biographical connection makes them useful for the \POSTECH--\KAIST{} Science
    War. International results already discussed above are not duplicated merely
    because a Korean mathematician or contestant was involved.

    \vspace{1em}

    {\small
    \renewcommand{\arraystretch}{1.18}
    \begin{longtable}{@{}L{0.09\textwidth}L{0.24\textwidth}L{0.27\textwidth}L{0.32\textwidth}@{}}
        \toprule
        Year & Person or event & Korean connection & Principal Quiz anchors \\
        \midrule
        \endfirsthead

        \toprule
        Year & Person or event & Korean connection & Principal Quiz anchors \\
        \midrule
        \endhead

        \midrule
        \multicolumn{4}{r}{Continued on the next page.} \\
        \endfoot

        \bottomrule
        \endlastfoot

        2025
        & Sug Woo Shin received the Samsung Ho-Am Prize in Physics and Mathematics
        & Distinguished professor at KIAS; professor at the University of California, Berkeley
        & Number theory, automorphic forms, the Langlands program. \\

        2025
        & Jae Kyoung Kim delivered a plenary lecture at the SIAM Annual Meeting
        & Chief investigator at the IBS Biomedical Mathematics Group
        & First Korean plenary speaker at a SIAM Annual Meeting; mathematical and
          computational biology. \\

        2025
        & Nam-Gyu Kang received the Korean Mathematical Society DI Prize
        & Professor at KIAS
        & Probability, complex analysis, conformal field theory. \\

        2025
        & Inhyeok Choi received the Sangsan Prize and joined KIAS
        & KIAS fellow and professor
        & Geometric group theory, random walks, hyperbolic groups. \\

        2025
        & Jae Choon Cha received the Korea Science Award
        & \POSTECH{} professor; Ph.D. from \KAIST
        & Geometric topology; Disk Embedding Theory in dimension four;
          transfinite invariants resolving a long-standing problem of John Milnor. \\

        2025
        & Cheolwoo Park received the Choi Seok-jeong Award
        & Professor at \KAIST
        & Statistical machine learning and statistical methodology. \\

        2026
        & Sung-Jin Oh received the Samsung Ho-Am Prize in Physics and Mathematics
        & \KAIST{} alumnus; professor at KIAS and associate professor at the University of California, Berkeley
        & Nonlinear hyperbolic PDE, general relativity, instability inside black holes. \\

        2026
        & Bo-Hae Im and JungHwan Park received Korean Mathematical Society Paper Awards
        & Professors at \KAIST
        & Positive-characteristic multiple zeta values and the Zagier--Hoffman
          conjectures; knot concordance and the $(2,1)$-cable of the figure-eight knot. \\

        2026
        & Bo-Hae Im, Hojin Kim, and Tuan Ngo Dac posted a preprint on Thakur's conjecture
        & Bo-Hae Im is a professor at \KAIST
        & Positive-characteristic multiple zeta values; independence through weight
          $2q$ and a unique relation at weight $2q+1$. \\

        2026
        & \POSTECH{} Arithmetic Geometry Lab selected for the Basic Research Laboratory Program
        & First collaborative arithmetic-geometry research group organized in the Republic of Korea; Sungmun Cho is principal investigator
        & Langlands program, Diophantine geometry, arithmetic geometry. \\

        2026
        & Kyeongsu Choi received the Science Prize at the POSCO TJ Park Prize
        & Professor at KIAS
        & Mean-curvature flow, Gauss-curvature flow, geometric PDE. \\

        2026
        & IMU Executive Committee met in Seoul
        & Hosted by KIAS on March 13--14
        & International Mathematical Union and Korean mathematical institutions. \\

        2026
        & Inaugural ASIACOMB conference in Daejeon, August 24--28
        & Held at the Daejeon Convention Center; organized as a new Asian
          combinatorics forum with IBS participation
        & Biennial combinatorics conference; plenary speakers include Jaehoon Kim,
          Alexander Postnikov, Van Vu, and Joshua Zahl. \\

        2025--2026
        & Jineon Baek and the moving sofa problem
        & \POSTECH{} alumnus; June E Huh Fellow at KIAS
        & Gerver's sofa, optimal area approximately $2.21953166$, preprint status. \\
    \end{longtable}
    }

    \begin{fact}[Jae Choon Cha and the 2025 Korea Science Award]
        Jae Choon Cha of \POSTECH{} was recognized for work in geometric
        topology. Two especially memorable anchors are his Disk Embedding Theory
        for four-dimensional topological manifolds and transfinite invariants of
        three-manifolds, the latter leading to a resolution of a problem posed by
        John Milnor that had remained open for more than sixty years. Cha received
        his Ph.D. from \KAIST, giving this entry a direct connection to both
        universities in the Science War.
    \end{fact}

    \begin{fact}[Two \KAIST{} faculty members and the 2026 KMS Paper Awards]
        The two \KAIST{} recipients have sharply different mathematical anchors.
        \begin{itemize}
            \item Bo-Hae Im was recognized for work determining algebraic
                  relations among multiple and alternating multiple zeta values
                  in positive characteristic and completely establishing the
                  positive-characteristic Zagier--Hoffman conjectures.
            \item JungHwan Park was recognized for work showing that the
                  $(2,1)$-cable of the figure-eight knot is not smoothly slice,
                  resolving a problem that had remained open for about forty
                  years and connecting directly to four-dimensional topology and
                  the slice--ribbon conjecture.
        \end{itemize}
    \end{fact}

    \begin{fact}[A 2026 preprint on Thakur's conjecture]
        Bo-Hae Im, Hojin Kim, and Tuan Ngo Dac posted the preprint
        \textit{The threshold for linear independence of multiple zeta values in
        positive characteristic}. For a finite field $\mathbb{F}_q$, they prove linear
        independence through weight $2q$, while at weight $2q+1$ there is a
        unique explicit $\mathbb{F}_q$-linear relation. Thus the preprint gives the first
        counterexample to Thakur's 2009 conjecture and identifies the precise
        threshold at which the conjectured independence fails. Because this work
        remained a preprint at the August 28, 2026 preparation cutoff, its
        status should be
        remembered together with the result.
    \end{fact}

    \begin{fact}[Two consecutive Samsung Ho-Am mathematics laureates]
        The Physics and Mathematics laureates in the two-year window are
        \[
            \begin{aligned}
                2025&:\quad \text{Sug Woo Shin}
                &&\longleftrightarrow\quad \text{Langlands program},\\
                2026&:\quad \text{Sung-Jin Oh}
                &&\longleftrightarrow\quad \text{nonlinear hyperbolic PDE and black holes}.
            \end{aligned}
        \]
        The prize category combines physics and mathematics, so the laureate's
        own discipline must be identified from the citation rather than from the
        category name alone.
    \end{fact}

    \begin{fact}[ASIACOMB 2026 in Daejeon]
        At the August 28 preparation cutoff, the inaugural ASIACOMB conference
        was taking place at the Daejeon Convention Center from August 24 through
        August 28. ASIACOMB is planned as a biennial conference in Asia devoted
        to combinatorics and international collaboration. Its plenary speakers
        include Jaehoon Kim of \KAIST, together with researchers such as
        Alexander Postnikov, Van Vu, and Joshua Zahl. Because the final day had
        not yet concluded at the cutoff, no unannounced prize result is recorded
        here.
    \end{fact}

    \begin{fact}[Korean competition results belong to the competition chapter]
        The Republic of Korea placed third at IMO 2025 and sixth at IMO 2026;
        Hyeonjun Lee's perfect 2026 score is recorded under Mathematical
        Competitions. These results are not repeated here, preserving one primary home for
        each event.
    \end{fact}

    \begin{strategy}
        Questions connected to the Republic of Korea are often easiest when three levels are kept
        separate:
        \[
            \text{nationality or origin},\qquad
            \text{current institution},\qquad
            \text{institution connected to the event}.
        \]
        A Korean or Korean-American mathematician may work abroad, hold a joint
        appointment, or receive a Korean prize for research performed in an
        international setting.
    \end{strategy}
